% =============================================================
%  附录 B  表与格
% =============================================================
\Chapter{附录 B\quad 表与格（随时复印）}

\begin{unitinfo}
本附录共 6 页，建议：\\
\textbf{完整口诀表}印 2 份（一份贴墙上，一份夹在手册里）；\\
\textbf{填空版口诀表}印 3 份（第七单元用，写坏了可以重来）；\\
\textbf{百数表}印 3 份（第四、七单元圈倍数用）；\\
\textbf{数位表与方格纸}各印 5 份（第一、二、三单元反复用）。
\end{unitinfo}

\clearpage
\subsection*{B1\quad 乘法口诀表（完整版·小数在前）}

\begin{center}
\begin{tabular}{@{}p{1.7cm}p{13.6cm}@{}}
\toprule
\textbf{1 的口诀} & 一一得一 \\
\midrule
\textbf{2 的口诀} & 一二得二\qquad 二二得四 \\
\midrule
\textbf{3 的口诀} & 一三得三\qquad 二三得六\qquad 三三得九 \\
\midrule
\textbf{4 的口诀} & 一四得四\qquad 二四得八\qquad 三四十二\qquad 四四十六 \\
\midrule
\textbf{5 的口诀} & 一五得五\qquad 二五一十\qquad 三五十五\qquad 四五二十\qquad 五五二十五 \\
\midrule
\textbf{6 的口诀} & 一六得六\qquad 二六十二\qquad 三六十八\qquad 四六二十四\qquad
                   五六三十\qquad 六六三十六 \\
\midrule
\textbf{7 的口诀} & 一七得七\qquad 二七十四\qquad 三七二十一\qquad 四七二十八\qquad
                   五七三十五\qquad 六七四十二\qquad 七七四十九 \\
\midrule
\textbf{8 的口诀} & 一八得八\qquad 二八十六\qquad 三八二十四\qquad 四八三十二\qquad
                   五八四十\qquad 六八四十八\qquad 七八五十六\qquad 八八六十四 \\
\midrule
\textbf{9 的口诀} & 一九得九\qquad 二九十八\qquad 三九二十七\qquad 四九三十六\qquad
                   五九四十五\qquad 六九五十四\qquad 七九六十三\qquad 八九七十二\qquad
                   九九八十一 \\
\bottomrule
\end{tabular}
\end{center}

\vspace{6pt}
\noindent\textbf{共 45 句。}检查一下：
$1+2+3+4+5+6+7+8+9=45$。

\begin{tipbox}
\textbf{怎么用这张表复习}：\\
① 从上往下，一行一行念（横着会）；\\
② 从左上往右下斜着念：一一得一、二二得四、三三得九……九九八十一（平方数锚点）；\\
③ 只念第 7 列：一七得七、二七十四……（竖着会）；\\
④ 豆角爸爸报前半句，小豆角报后半句，\textbf{跳着来}。
\end{tipbox}

\clearpage
\subsection*{B2\quad 乘法口诀填空表（练习用）}

\begin{center}
\begin{tabular}{@{}p{1.7cm}p{13.6cm}@{}}
\toprule
\textbf{1 的口诀} & 一一 \blank[1.4cm] \\
\midrule
\textbf{2 的口诀} & 一二 \blank[1.4cm]\qquad 二二 \blank[1.4cm] \\
\midrule
\textbf{3 的口诀} & 一三 \blank[1.4cm]\qquad 二三 \blank[1.4cm]\qquad 三三 \blank[1.4cm] \\
\midrule
\textbf{4 的口诀} & 一四得四\qquad 二四 \blank[1.4cm]\qquad 三四 \blank[1.4cm]\qquad 四四 \blank[1.4cm] \\
\midrule
\textbf{5 的口诀} & 一五 \blank[1.4cm]\qquad 二五 \blank[1.4cm]\qquad 三五 \blank[1.4cm]\qquad
                   四五 \blank[1.4cm]\qquad 五五 \blank[1.4cm] \\
\midrule
\textbf{6 的口诀} & 一六 \blank[1.4cm]\qquad 二六 \blank[1.4cm]\qquad 三六 \blank[1.4cm]\qquad
                   四六 \blank[1.4cm]\qquad 五六 \blank[1.4cm]\qquad 六六 \blank[1.4cm] \\
\midrule
\textbf{7 的口诀} & 一七 \blank[1.4cm]\qquad 二七 \blank[1.4cm]\qquad 三七 \blank[1.4cm]\qquad
                   四七 \blank[1.4cm]\qquad 五七 \blank[1.4cm]\qquad 六七 \blank[1.4cm]\qquad
                   七七 \blank[1.4cm] \\
\midrule
\textbf{8 的口诀} & 一八 \blank[1.4cm]\qquad 二八 \blank[1.4cm]\qquad 三八 \blank[1.4cm]\qquad
                   四八 \blank[1.4cm]\qquad 五八 \blank[1.4cm]\qquad 六八 \blank[1.4cm]\qquad
                   七八 \blank[1.4cm]\qquad 八八 \blank[1.4cm] \\
\midrule
\textbf{9 的口诀} & 一九 \blank[1.4cm]\qquad 二九 \blank[1.4cm]\qquad 三九 \blank[1.4cm]\qquad
                   四九 \blank[1.4cm]\qquad 五九 \blank[1.4cm]\qquad 六九 \blank[1.4cm]\qquad
                   七九 \blank[1.4cm]\qquad 八九 \blank[1.4cm]\qquad 九九 \blank[1.4cm] \\
\bottomrule
\end{tabular}
\end{center}

\noindent\textbf{批改速查}：得数依次是 \ 1；2、4；3、6、9；4、8、12、16；
5、10、15、20、25；6、12、18、24、30、36；7、14、21、28、35、42、49；
8、16、24、32、40、48、56、64；9、18、27、36、45、54、63、72、81。

\clearpage
\subsection*{B3\quad 百数表（圈倍数用）}

\begin{center}
\setlength{\tabcolsep}{2pt}
\def\arraystretch{1.55}
\begin{tabular}{|>{\centering\arraybackslash}p{1.42cm}|*{9}{>{\centering\arraybackslash}p{1.42cm}|}}
\hline
1 & 2 & 3 & 4 & 5 & 6 & 7 & 8 & 9 & 10 \\ \hline
11 & 12 & 13 & 14 & 15 & 16 & 17 & 18 & 19 & 20 \\ \hline
21 & 22 & 23 & 24 & 25 & 26 & 27 & 28 & 29 & 30 \\ \hline
31 & 32 & 33 & 34 & 35 & 36 & 37 & 38 & 39 & 40 \\ \hline
41 & 42 & 43 & 44 & 45 & 46 & 47 & 48 & 49 & 50 \\ \hline
51 & 52 & 53 & 54 & 55 & 56 & 57 & 58 & 59 & 60 \\ \hline
61 & 62 & 63 & 64 & 65 & 66 & 67 & 68 & 69 & 70 \\ \hline
71 & 72 & 73 & 74 & 75 & 76 & 77 & 78 & 79 & 80 \\ \hline
81 & 82 & 83 & 84 & 85 & 86 & 87 & 88 & 89 & 90 \\ \hline
91 & 92 & 93 & 94 & 95 & 96 & 97 & 98 & 99 & 100 \\ \hline
\end{tabular}
\end{center}

\begin{tipbox}
\textbf{怎么用（每一样都做一次）}：\\
① 用蓝笔圈出 \textbf{2 的倍数}（2、4、6……），看看形成什么形状；\\
② 用红笔圈出 \textbf{5 的倍数}，看看它们排在哪一列；\\
③ 用铅笔圈出 \textbf{3 的倍数}，数一数有几条斜线；\\
④ 用黄笔涂出 \textbf{9 的倍数}，看看斜线上的数字各位相加是几。\\
⑤ 把 \textbf{平方数}（1、4、9、16……81）用方框框出来。
\end{tipbox}

\clearpage
\subsection*{B4\quad 数位表（写竖式、读数用）}

\begin{center}
\renewcommand{\arraystretch}{2.6}
\begin{tabular}{|>{\centering\arraybackslash}p{3.2cm}|*{3}{>{\centering\arraybackslash}p{3.2cm}|}}
\hline
\textbf{百位} & \textbf{十位} & \textbf{个位} & \\ \hline
 & & & \\ \hline
 & & & \\ \hline
 & & & \\ \hline
 & & & \\ \hline
 & & & \\ \hline
\end{tabular}
\end{center}

\vspace{4pt}
\begin{center}
\renewcommand{\arraystretch}{2.2}
\begin{tabular}{|>{\centering\arraybackslash}p{2.4cm}|*{3}{>{\centering\arraybackslash}p{2.4cm}|}>{\centering\arraybackslash}p{2.4cm}|}
\hline
\textbf{千位} & \textbf{百位} & \textbf{十位} & \textbf{个位} & \textbf{读作} \\ \hline
 & & & & \\ \hline
 & & & & \\ \hline
 & & & & \\ \hline
\end{tabular}
\end{center}

\begin{tipbox}
\textbf{怎么用}：\\
① 豆角爸爸说一个数（如“三百零五”），小豆角在对应数位格子里写数字；\\
② 反过来，格子里写数字，让小豆角读出来；\\
③ 写竖式时，\textbf{一个格子写一个数字}——这就是“数位对齐”的物理保证。
\end{tipbox}

\clearpage
\subsection*{B5\quad 方格纸（1 厘米方格，画线段与点子图用）}

\begin{center}
\begin{tikzpicture}
  \draw[step=1cm,gray!45,very thin] (0,0) grid (16,9);
  \draw[mainc,thick] (0,0) rectangle (16,9);
\end{tikzpicture}
\end{center}

\vspace{8pt}
\begin{center}
\begin{tikzpicture}
  \draw[step=1cm,gray!45,very thin] (0,0) grid (16,9);
  \draw[mainc,thick] (0,0) rectangle (16,9);
\end{tikzpicture}
\end{center}

\begin{tipbox}
\textbf{怎么用}：\\
① 画线段：起点落在格点上，长度就等于格数（画 5 厘米就数 5 格）；\\
② 点子图：在每个格点上画一个点，就是“几行几列”的点子图；\\
③ 长方形面积（预习）：圈出一个长方形，数里面的格子数，
就是“每行几个 $\times$ 几行”。
\end{tipbox}

\clearpage
\subsection*{B6\quad 竖式练习格（第一单元专用）}

\noindent 每格一个数位，写竖式不会错位。建议印 3 份。

\vspace{6pt}
\begin{center}
\newcommand{\vsframe}{%
\begin{tikzpicture}
  \draw[mainc!60,thick] (0,0) rectangle (4.4,3.6);
  \draw[gray!40] (0,1.2) -- (4.4,1.2);
  \draw[gray!40] (0,2.4) -- (4.4,2.4);
  \draw[gray!40] (1.1,0) -- (1.1,3.6);
  \draw[gray!40] (2.2,0) -- (2.2,3.6);
  \draw[gray!40] (3.3,0) -- (3.3,3.6);
\end{tikzpicture}}
\vsframe\hspace{0.6cm}\vsframe\hspace{0.6cm}\vsframe

\vspace{0.6cm}
\vsframe\hspace{0.6cm}\vsframe\hspace{0.6cm}\vsframe

\vspace{0.6cm}
\vsframe\hspace{0.6cm}\vsframe\hspace{0.6cm}\vsframe

\vspace{0.6cm}
\vsframe\hspace{0.6cm}\vsframe\hspace{0.6cm}\vsframe
\end{center}

\begin{tipbox}
\textbf{怎么用}：最右一列写个位，中间写十位，最左写百位。
\textbf{进位的小 1 写在十位格子的上边}，\textbf{退位的小点画在十位数字的头顶}。
12 个格子，够做 12 道竖式。
\end{tipbox}
